% Early vs late: what late pays for, the threshold, cache, and choosing per chunk.

Late materialization reads the predicate column in full, and of every other column only the
lines that hold a selected row. Early streams every referenced column in full. This section
measures why late's smaller line count does not always win, turns that into a threshold, and
shows how the answer changes when the data is in cache.

\subsection{What late pays for}
\label{sec:mat:cost}

\begin{table}[t]
\centering\small
\caption{Lines late reads from one gathered column, per 8{,}000 rows (1{,}000 lines of 8
values). Lines $= 1000 \cdot f(s)$.}
\label{tab:linecount}
\begin{tabular}{rrrr}
\toprule
Selected & Values wanted & Lines read & Values brought in \\
\midrule
1\% & 80 & 77 & 616 \\
5\% & 400 & 337 & 2{,}696 \\
10\% & 800 & 570 & 4{,}560 \\
\bottomrule
\end{tabular}
\end{table}

\paragraph{Late reads lines, not values.} \Cref{tab:linecount} shows how fast $f(s)$ grows. At
1\% late touches 8\% of a column's lines, at 5\% a third, at 10\% more than half, while it wants
only one value from each. Two further costs, which the line count of \cref{tab:linemodel}
leaves out, decide where late stops winning. We separated them by timing late and early on DSM
for $k = 1$ to 16 at 0.01--10\% (\code{results/late\_breakdown/}).
\begin{itemize}
\item \textbf{A fixed first pass.} Building the selection vector with scalar code costs about 2
cycles per row, against about 1.2 for streaming 8 bytes from DRAM. Late with $k{=}1$ at
0.01\% takes 2.05\ms{} for 4M rows. With the AVX2 selection vector (\cref{sec:mat:cache}) it
takes 1.21\ms, the time of one bare column read, so the scalar pass is CPU-bound and costs
about 0.7 column scans more than reading $c_0$.
\item \textbf{A gathered line costs more than a streamed one}, 2.2--3.2 times per touched line.
A gathered column takes 0.34\ms{} at 1\% (8\% of its lines), 0.66\ms{} at 2.4\% (18\%),
1.17\ms{} at 5\% (34\%) and 1.97\ms{} at 10\% (57\%), against 1.27--1.63\ms{} for a whole
column in early. Gathered lines are sparse and irregular, so the prefetcher cannot run ahead of
them, and each is a separate DRAM access limited by latency and the number of misses in flight.
At 5\%, gathering a third of a column (1.28\ms{} with $k{=}2$) already takes longer than
reading all of it (1.21\ms).
\end{itemize}
The first pass is paid once per query, and every column late skips pays some of it back. That
is why the crossover grows with $k$: 1.6\% for $k{=}2$, 2.4\% for $k{=}4$, 7.4\% for $k{=}16$
(\cref{tab:crossovers}).

\subsection{A threshold from the line size}
\label{sec:mat:threshold}

In units of one sequential column scan per row, early costs $1 + g$ and late
$1 + \sigma + g\,f(s)/\varphi$, where $g$ is the number of attributes late gathers, $1/\varphi$
the cost of a gathered line relative to a streamed one, and $\sigma$ late's fixed cost. Late
wins while
\[
  f(s) < \varphi\,(1-\sigma/g)
  \quad\Longleftrightarrow\quad
  s < 1-\bigl(1-\varphi(1-\sigma/g)\bigr)^{1/v}.
\]
The machine enters through $v = \text{line size}/\text{value width}$. Fitted to the DSM
crossovers from DRAM, $\varphi = 0.55$ and $\sigma = 2$, and the rule predicts 2.5\% for $g=3$
and 7.8\% for $g=15$, against 2.4\% and 7.4\% measured. The two constants are a joint fit, not
the costs measured one at a time (0.7 scans and 2.2--3.2$\times$): the model charges early a
constant cost per column and assumes a constant $\varphi$, which the measurements do not, and
with the measured values it would predict 4.4\% and 5.5\%. The rule also shows two things.
Late can never win when $g \le \sigma$, i.e.\ when it skips at most two attributes, and for
$g=1$ late is indeed at most 18\% faster. And packing more values into a line lowers the
threshold: a 10-bit column holds 51 values per line, so $f(1\%)$ is already 0.40
(\cref{sec:enc:compressed}).

\subsection{In cache, instructions per row decide}
\label{sec:mat:cache}

\begin{table}[t]
\centering
\caption{Selectivity at which late stops beating early with the table in L3 (4\,MB, compare
\cref{tab:crossovers}).}
\label{tab:crossovers_cache}
\gen{selectivity_cache_crossovers}
\end{table}

\begin{table}[t]
\centering\small
\caption{Where late stops beating early on DSM ($k{=}4$), with the scalar and the AVX2
selection vector.}
\label{tab:simdcross}
\gen{simd_crossovers}
\end{table}

With the table in L3 (\cref{tab:crossovers_cache}), the crossovers move apart in opposite
directions. Projections favour late much longer: 20\% instead of 2.4\% for $k{=}4$ on DSM, 33\%
instead of 7.4\% for $k{=}16$, because a gathered line now costs an L3 hit instead of a DRAM
access, while early still pays the CPU time of copying every row. SUM goes the other way: late
never wins on DSM or PAX.

The reason is that cache shrinks early's cost far more than late's. For SUM of 4 DSM columns,
early's cost is bytes: 4.4 cycles per row from DRAM, 1.2 in L3 (0.30\,ns), where the reads are
about three times faster and early's masked sum (\code{total += row\_sum \& mask}: no branch,
no store, the same work for every row) is vectorized by the compiler. Late's first pass is
instructions: about 2.1 cycles per row from DRAM and 1.9 in L3 (0.46\,ns). In L3, late spends
more on building the selection vector alone than early spends on the whole query, so it loses
even at 0.01\%, where it skips almost every line.

\paragraph{A hand-written SIMD selection vector.} The compiler cannot vectorize the selection
loop, although the build allows AVX2 (\code{-march=native}). Its output is a packed list:
where row $i$ goes depends on how many earlier rows passed (a stream compaction), so the lanes
of a register cannot be processed independently. Our AVX2 kernel compares 8 values with two
4-lane 64-bit compares (\code{vpcmpgtq}, signed only, so the sign bit is flipped to make the
range test unsigned), turns the result into an 8-bit mask (\code{movemask}), loads the
positions of the set bits from a 256-entry table (8\,KB, stays in L1), adds the vector offset,
stores all 8 and advances the output index by \code{popcount}(mask). AVX2 has no compress
instruction; AVX-512's \code{vpcompressd} would replace the table, but Zen~3 has no AVX-512.
\begin{itemize}
\item In cache the first pass drops from about 2 to 0.9 cycles per row (0.46 to 0.20\,ns in
L3, 0.51 to 0.24 in L2), now below early's 1.2, and late SUM wins up to 3.4\% in L3 and 5.8\%
in L2 (\cref{tab:simdcross}).
\item From DRAM, late is memory-bound and gains about a third only at the lowest
selectivities (the first pass falls from 2.05 to 1.21\ms); the crossovers move from 1.2--2\% to
about 3\%.
\item Early does not use a selection vector and is unchanged.
\end{itemize}

\subsection{Choosing per chunk}
\label{sec:mat:auto}

\begin{table}[t]
\centering\small
\caption{Auto against the better of early and late at each of 10 selectivities, with the table
in L3: geometric mean / worst point over NSM, DSM and PAX (1.00 is always right).}
\label{tab:autocache}
\begin{tabular}{lccc}
\toprule
Query & DRAM constants & in-cache constants & measured (\cref{sec:tried}) \\
\midrule
SUM, $k{=}4$ & 1.19 / 1.68 & 1.20 / 1.55 & 1.09 / 1.16 \\
\code{SELECT} $k{=}4$ & 1.06 / 1.35 & 1.05 / 1.20 & 1.06 / 1.12 \\
\code{SELECT *} & 1.13 / 2.58 & \textbf{1.02 / 1.30} & 1.06 / 1.10 \\
\bottomrule
\end{tabular}
\end{table}

\emph{auto} applies the threshold to every span of at least 1024 rows, using the selectivity of
the previous span. Spans are at least 1024 rows because a 128-row PAX page at 1\% holds 0 to 4
qualifying rows, too few to estimate a selectivity from.
\begin{itemize}
\item \textbf{Where the constants were fitted it is enough.} From DRAM on uniform data, auto is
within 1\% of the better fixed strategy on NSM and DSM at every selectivity (geometric mean);
always running late would cost 21\% on DSM, always early 42\%.
\item \textbf{On sorted data it beats both fixed strategies}, by 10\% on DSM on average and by
up to 24\% (\code{SELECT *} at 10\%). The qualifying rows form one contiguous stretch, so
late is right outside it and early inside, and choosing per span gets both.
\item \textbf{The constants depend on where the data lives.} In L3, the DRAM constants pick
early for \code{SELECT *} at 10--20\%, 2.3--2.6$\times$ too slow. Since the engine knows the
table's size, \code{--auto-cache} switches to a second pair, fitted on the in-cache runs
($\varphi = 1$, $\sigma = 2$: a gathered line costs about a streamed one), whenever the table
fits in L3. That brings \code{SELECT *} to 1.02 on average and 1.30 at worst
(\cref{tab:autocache}).
\item \textbf{The remaining misses need more constants.} In cache, early SUM vectorizes and
early \code{SELECT} does not, so one pair cannot fit both, and SUM stays up to 1.55$\times$ off.
On PAX from DRAM, whose early strategy is slower than DSM's, \code{SELECT *} is up to
2.4$\times$ off. A threshold per query class and layout would fix both, at the price of more
fitted constants.
\end{itemize}
